% =========================================================
% Proof Stub: Dense Subsequence Criterion
% Source: volume-iii/analysis/sequences/notes/sequences/notes-subsequence-toolkit.tex
% Mode: standard
% =========================================================

\clearpage
\phantomsection
\hypertarget{proof-RA-ABB-T-dense-subseq-criterion}{}

\subsubsection[Dense Subsequence Criterion]{Proof --- Dense Subsequence Criterion}

\bigskip

\noindent
\textbf{Source.}
See \texttt{notes-subsequence-toolkit.tex}.

\vspace{0.75em}

\noindent
\textbf{Goal.}
If every neighborhood of $L$ contains infinitely many terms of $(a_n)$, then $(a_n)$ has a subsequence converging to $L$.

\vspace{0.75em}

\noindent
\textbf{Logical form.}
\[
  \forall \varepsilon > 0,\; \{n : |a_n - L| < \varepsilon\} \text{ infinite} \Rightarrow \exists \text{ subseq } (a_{n_k}) \to L.
\]

\vspace{0.75em}

\noindent
\textbf{Proof strategy.}
Recursively select indices $n_1 < n_2 < \cdots$ with $|a_{n_k} - L| < 1/k$. The infinitely-many-terms hypothesis guarantees each selection can be made with $n_k > n_{k-1}$.

\vspace{0.75em}

\noindent
\textbf{Proof.}
\begin{proof}
\Claim If every neighborhood of $L$ contains infinitely many terms of $(a_n)$, then $(a_n)$ has a subsequence converging to $L$.

\Given 

\Goal 

\Strategy Recursively select indices $n_1 < n_2 < \cdots$ with $|a_{n_k} - L| < 1/k$. The infinitely-many-terms hypothesis guarantees each selection can be made with $n_k > n_{k-1}$.

\medskip

% --- for k = 1: use ε = 1, pick n_1 with |a_{n_1} - L| < 1 ---
% --- for k+1: use ε = 1/(k+1); infinitely many n satisfy this; pick n_{k+1} > n_k ---
% --- by construction n_1 < n_2 < ... and |a_{n_k} - L| < 1/k → 0 ---
% --- therefore a_{n_k} → L ---

\end{proof}

\begin{remark*}[Proof shape]
Build the subsequence greedily: at stage $k$ use $\varepsilon = 1/k$ and pick the next available index.
\end{remark*}

\begin{remark*}[Dependencies]
\begin{itemize}
  \item Infinitely many terms in each $\varepsilon$-ball (hypothesis).
  \item Squeeze theorem (or direct: $1/k \to 0$).
  \item Definition of subsequence.
\end{itemize}
\end{remark*}
